% !TEX program = pdflatex
% Amazon Research Awards Proposal
% Compile: pdflatex main.tex
%
\documentclass[11pt]{ara-proposal}

\title{Project Title}

\piName{<PI name>}
\piTitle{<title>}
\piDepartment{<department name>}
\piUniversity{<university name>}

\copiName{<Co-PI name>}
\copiTitle{<title>}
\copiDepartment{<department name>}
\copiUniversity{<university name>}

\cashFunding{<amount>}
\awsCredits{<amount>}
\amazonContact{<Amazon contact name>}{<email>}

\keywords{keyword 1, keyword 2}

\begin{document}

\makeAraTitle

% ==============================
% ABSTRACT
% ==============================
\begin{abstract}
This is the abstract section. Replace this text with a concise summary of your
project, its objectives, methods, and expected impact.
\end{abstract}

% ==============================
% INTRODUCTION
% ==============================
\section{Introduction}
Significance of the research and prior work. 
Provide context, motivation, and a brief overview of the current state of the art.

% ==============================
% METHODS
% ==============================
\section{Methods}
Your technical approach.

% ==============================
% EXPECTED RESULTS
% ==============================
\section{Expected Results}
Please include milestones with timeline estimates, such as for datasets, code releases, technical reports,
publications, applications, presentations, etc

\section{Fund Needed}
Please provide justification supporting the designated amounts for cash funding and AWS Promotional Credits. For cash funding, please do not include indirect expenses which are not incurred solely for your project. For AWS Promotional Credits, please list the AWS products you plan to use.

\section{Additional Information}
\emph{What hard problem does your research address, and why has it resisted solution at scale? How does your approach differ from existing methods?}

Bla bla bla

\emph{What would make this research impactful rather than incremental? What are the key milestones between your research and practical adoption?}

Bla bla bla

\emph{Why is now the right time to pursue this work? What recent developments or enabling factors make it feasible?}

Bla bla bla

\emph{What are the key risks and how will you mitigate them? What safety considerations or potential dual-use concerns does your research raise, and how will you address them?}

Bla bla bla

\emph{How might your research advance the field broadly? What are potential applications of your work that could benefit society, including potential applications at organizations like Amazon?}

Bla bla bla

% ==============================
\appendix

\newappendix{Appendix A\,--\,References \textnormal{\small(does not count toward page limit)}}
Limit two pages.

\begin{thebibliography}{99}
\bibitem{example1} Author, ``Title of the work,'' \emph{Venue}, Year.
\end{thebibliography}

\newappendix{Appendix B\,--\,CV of PI \textnormal{\small(does not count toward page limit)}}
One-page CV of PI (professional title, experience, etc.), only to include the 5--6
most relevant papers to this proposal. One-page CV of Co-PI (optional).

\newappendix{Appendix C\,--\,Previously Funded Project Summary \textnormal{\small(does not count toward page limit)}}
If you or someone on your project team has received cash funding or AWS Promotional
Credits directly or indirectly from Amazon within the past 3 years, please answer
the questions below. (Include funding from Amazon through your institution or from
the NSF-Amazon Fairness in AI funding grant.)

\begin{description}
  \item[Project title:]~\\
  \item[Funding amount and time:]~\\
  \item[Funding source:]~\\
  \item[Who was this funding for?]~\\
  \item[If AWS Promotional Credits were provided, how much is left?]~\\
  \item[Summary:] Please provide a brief summary of the results, papers, and
    open-source software enabled.
\end{description}

\end{document}
